\section{Complex algebraic orbit closures}

\begin{lemma}[Intrinsic polynomials in coordinates]
\label{lem:complex-closure-coordinates}
\uses{lem:balanced-intrinsic-polynomials,lem:balanced-invariant-sufficient}
\lean{BostConjecture.InvariantTheory.polynomialEval_equivMvPolynomial_symm,BostConjecture.InvariantTheory.polynomialEval_eq_eval_equivMvPolynomial,BostConjecture.InvariantTheory.polynomialZariskiClosure_eq_preimage_zeroLocus,BostConjecture.InvariantTheory.isClosed_zeroLocus,BostConjecture.InvariantTheory.continuous_polynomialEval,BostConjecture.InvariantTheory.isClosed_polynomialZariskiClosure,BostConjecture.InvariantTheory.closure_subset_polynomialZariskiClosure}\leanok
Let $E$ be a finite-dimensional vector space over a field $K$, with basis
$(b_i)_{i\in T}$ and coordinate isomorphism $c:E\to K^T$. The canonical
algebra isomorphism
\[
 \operatorname{Sym}(E^*)\longrightarrow K[X_i\mid i\in T],
 \qquad b_i^*\longmapsto X_i,
\]
intertwines evaluation at $x$ and at $c(x)$. Consequently, for every
$S\subseteq E$, its intrinsic polynomial Zariski closure is
$c^{-1}(V(I(c(S))))$, where $I$ denotes the vanishing ideal.
If $K=\mathbb C$ and $E$ carries a complex norm, this closure is
Euclidean closed and contains the Euclidean closure of $S$.
\end{lemma}
\begin{proof}
Both evaluation maps are algebra homomorphisms and agree on each generator
$b_i^*$, whose value is the $i$th coordinate of $x$. Thus they agree on
the whole symmetric algebra. The isomorphism is surjective, so testing
all intrinsic polynomials vanishing on $S$ is equivalent to testing all
coordinate polynomials vanishing on $c(S)$. This gives the asserted
zero-locus identity. In the complex case the coordinate map is continuous
because every linear map on a finite-dimensional normed space is continuous.
Each coordinate polynomial is continuous, so its zero set is closed.
Arbitrary intersections of these zero sets, and their preimages under $c$,
are closed. The intrinsic polynomial closure contains $S$ by definition,
and therefore contains its Euclidean closure.
\end{proof}

\begin{lemma}[Persistence of roots of monic polynomials]
\label{lem:complex-closure-root-persistence}
\lean{BostConjecture.InvariantTheory.exists_roots_tendsto}\leanok
Let $F_k\in\mathbb C[T]$ be monic of a fixed positive degree $r$, and
let $z\in\mathbb C$. If $F_k(z)\to0$, there are roots $z_k$ of $F_k$
such that $z_k\to z$.
\end{lemma}
\begin{proof}
Each $F_k$ has exactly $r$ roots counted with multiplicity, and at least
one root since $r>0$. Choose $z_k$ minimizing distance to $z$ among
these roots. The monic factorization gives
\[
 |z_k-z|^r\leq\prod_{w:F_k(w)=0}|z-w|=|F_k(z)|,
\]
where the product counts multiplicities. For each $\epsilon>0$, the
right side is eventually less than $\epsilon^r$, so
$|z_k-z|<\epsilon$ eventually. This is precisely $z_k\to z$;
no assumption on multiplicities is needed.
\end{proof}

\begin{lemma}[Uniform bounds for roots under specialization]
\label{lem:complex-closure-root-bounds}
\lean{BostConjecture.InvariantTheory.exists_bound_roots_specialization}\leanok
Let $F\in\mathbb C[X_1,\ldots,X_d][T]$ be monic, and let
$a_k\to a$ in $\mathbb C^d$. There is a constant $B$ such that every
root of every $F(a_k,T)$ has modulus at most $B$.
\end{lemma}
\begin{proof}
Writing $r=\deg_T F$, monicity ensures that every specialization has
degree $r$ and leading coefficient $1$. The Cauchy bound gives, for each
root $w$,
\[
 |w|<1+\max_{i<r}|F_i(a_k)|
       \leq 1+\sum_{i<r}|F_i(a_k)|.
\]
The right side is a convergent real sequence, since each coefficient is
a continuous polynomial, so it is bounded. For $r=0$ the monic polynomial
is $1$ and has no roots, and the assertion is vacuous.
\end{proof}

\begin{lemma}[Constructible images]
\label{lem:balanced-constructible-image}
\source{milne-algebraic-geometry:page-0201}\notready
A regular map of complex affine algebraic varieties sends constructible
sets to constructible sets.
\end{lemma}
\begin{proof}
We first prove that a dominant map of irreducible affine varieties has
image containing a nonempty open subset. Its map of coordinate rings
is an inclusion $A\subseteq B$ of finitely generated complex domains,
with $B$ finitely generated over $A$. Choose a maximal algebraically
independent subset $t_1,\ldots,t_r$ of a finite set of generators of
$B$ over $\operatorname{Frac}(A)$. The remaining generators are
algebraic over $\operatorname{Frac}(A)(t_1,\ldots,t_r)$. Clear the
finitely many denominators in their algebraic equations and invert
the product $q(t)$ of their leading coefficients. After also inverting
some $0\ne a\in A$, the algebra
$B[1/a,1/q]$ is finite integral over $A[1/a][t_1,\ldots,t_r,1/q]$.
Multiply $a$ by a nonzero numerator of one coefficient of $q$, so the specialization
$q_y$ is nonzero at every closed point $y\in D(a)$.
Since $\mathbb C$ is infinite, choose $t\in\mathbb C^r$ with
$q_y(t)\ne0$. Evaluation at $(y,t)$ gives a maximal ideal of the
latter algebra. Lying over supplies a maximal ideal of the former;
its residue field is $\mathbb C$ by the weak Nullstellensatz. Thus
$y$ lies in the image, proving the assertion about $D(a)$.

For the general assertion, use Noetherian induction on the closure of
the constructible set $C$. Splitting into finitely many irreducible
components and restricting to affine open subsets reduces to $C$ dense
in an irreducible affine variety $Y$. Expressing $C$ as a finite union
of locally closed subsets shows that one of those subsets contains a
nonempty open subset $V$ of $Y$: otherwise their closures could not
cover the irreducible $Y$. Choose a nonempty affine open inside $V$.
Apply the first paragraph to its map into the closure $Z$ of the
image. There is a nonempty open $O\subseteq Z$ contained in the
image of $C$. Then
$f(C)=O\cup f(C\cap f^{-1}(Z\setminus O))$.
The second term is constructible by induction, because its source
has proper closed closure in $Y$. This proves the assertion.
This is the affine form of \cite[Theorem 9.7]{milne-algebraic-geometry}.
\end{proof}

\begin{lemma}[Lifting roots of the actual minimal polynomial]
\label{lem:complex-closure-integral-lifting}\notready
Let $K$ be algebraically closed, let $R\subseteq B$ be an inclusion of
$K$-domains with $R$ integrally closed, and suppose $B$ is integral over
$R$ and finitely generated over $K$. For $u\in B$, let $F\in R[T]$ be
its monic minimal polynomial. If $\chi:R\to K$ is a $K$-algebra
homomorphism and $F^\chi(z)=0$, then there is a $K$-algebra
homomorphism $\psi:B\to K$ extending $\chi$ with $\psi(u)=z$.
\end{lemma}
\begin{proof}
The kernel of $R[T]\to B$, $T\mapsto u$, is $(F)$: divide by the
monic $F$ and use minimality over the fraction field to make the remainder
zero. Integrally closedness ensures that this monic minimal polynomial
has coefficients in $R$. Thus evaluation of coefficients by $\chi$ and
of $T$ by $z$ descends to a $K$-algebra homomorphism
$\phi:R[u]\to K$. Its kernel is maximal since $\phi$ contains all
constants in its image. The extension $R[u]\subseteq B$ is integral,
so lying over provides a maximal ideal $Q\subset B$ contracting to
$\ker\phi$. The field $B/Q$ is finitely generated as a $K$-algebra,
hence finite algebraic over $K$ by the weak Nullstellensatz and equal
to $K$ since $K$ is algebraically closed. Composing the quotient map
with this $K$-algebra identification gives $\psi$. For $a\in R[u]$,
$a-\phi(a)$ belongs to $Q$, so $\psi(a)=\phi(a)$. In particular
$\psi$ extends $\chi$ and sends $u$ to $z$.
\end{proof}

\begin{lemma}[Polynomial separation on a finite set]
\label{lem:complex-closure-finite-separator}\notready
Let $S$ be a finite subset of $K^T$, where $K$ is a field, and let
$p\in K^T\setminus S$. There is a polynomial $u\in K[X_i\mid i\in T]$
such that $u(p)=1$ and $u(q)=0$ for every $q\in S$.
\end{lemma}
\begin{proof}
For each $q\in S$, choose $j(q)\in T$ with $p_{j(q)}\ne q_{j(q)}$.
Set
\[
 u(X)=\prod_{q\in S}
 \frac{X_{j(q)}-q_{j(q)}}{p_{j(q)}-q_{j(q)}}.
\]
All denominators are nonzero constants. Every factor is one at $p$,
and the factor indexed by $q$ vanishes at $q$. The empty product is one,
so the assertion includes $S=\varnothing$ and empty coordinate arity.
\end{proof}

\begin{lemma}[Complex and Zariski constructible closures]
\label{lem:balanced-constructible-closure}
\uses{lem:complex-closure-root-persistence,lem:complex-closure-root-bounds,lem:complex-closure-coordinates,lem:complex-closure-integral-lifting,lem:complex-closure-finite-separator}
\source{milne-algebraic-geometry:page-0201}\notready
A constructible subset of a complex affine variety has the same
Euclidean and Zariski closures.
\end{lemma}
\begin{proof}
It suffices to prove that a nonempty Zariski open subset of an
irreducible affine variety $Y\subseteq\mathbb C^N$ is Euclidean dense:
decompose a constructible set into finitely many irreducible locally closed
pieces and take their closures. Let $B$ be the coordinate domain of $Y$.
Noether normalization gives a finite inclusion
$A=\mathbb C[t_1,\ldots,t_d]\subseteq B$, inducing
$\pi:Y\to\mathbb C^d$. We first show that $\pi$ is Euclidean open.

Every ambient coordinate satisfies a monic polynomial over $A$. Its values
on a fixed fiber belong to a finite root set, so each fiber is finite.
Fix $p\in Y$, put $a=\pi(p)$, and for each other point $q$ of this fiber
choose a coordinate $j(q)$ such that $p_{j(q)}\ne q_{j(q)}$. The regular function
\[
 u(x)=\prod_{q\in\pi^{-1}(a)\setminus\{p\}}
 \frac{x_{j(q)}-q_{j(q)}}{p_{j(q)}-q_{j(q)}}
\]
is $1$ at $p$ and $0$ at all other points of the fiber; an empty product
is $1$. Since $A$ is integrally closed and $u$ is integral over $A$, its
monic minimal polynomial over $\operatorname{Frac}(A)$ has coefficients in
$A$. Call it $F$, of degree $r>0$. Monic division and minimality show
that the kernel of $A[T]\to B$, $T\mapsto u$, is exactly $(F)$.
Thus $A[u]=A[T]/(F)$. Every complex solution $(a',z)$ of $F(a',z)=0$
defines a maximal ideal of $A[u]$. The extension $A[u]\subseteq B$ is
integral, so lying over and the weak Nullstellensatz give $q\in Y$
with $\pi(q)=a'$ and $u(q)=z$.

Let $a_k\to a$. The specialization $F(a_k,T)$ remains monic of degree
$r$. Among its roots choose $z_k$ nearest to $1$. Factorization over
$\mathbb C$ gives
\[
 |z_k-1|^r\leq |F(a_k,1)|\longrightarrow |F(a,1)|=0,
\]
so $z_k\to1$. Lift $(a_k,z_k)$ to $q_k\in Y$ as above. The
coefficients of each ambient coordinate's monic integral equation are
bounded along $(a_k)$. The Cauchy bound for polynomial roots therefore
bounds every coordinate of $(q_k)$. Every subsequence has a convergent
further subsequence. Its limit $q$ lies in $Y$, since polynomial zero
sets are Euclidean closed, and continuity gives $\pi(q)=a$ and $u(q)=1$.
The separating function forces $q=p$. Hence the whole sequence converges
to $p$: otherwise a subsequence outside a fixed neighborhood of $p$ would
have a further subsequence converging to $p$, a contradiction.
If the image of an open neighborhood $V$ of $p$ were not a neighborhood
of $a$, one could choose $a_k\notin\pi(V)$ tending to $a$; the lifts
just constructed would eventually belong to $V$. This contradiction
proves openness of $\pi$.

Now let $Z\subsetneq Y$ be closed. Choose a nonzero $b\in B$ vanishing
on $Z$. In a monic integral relation for $b$, remove the largest power
of $T$ dividing the relation. Cancellation is valid because $B$ is a
domain and $b\ne0$. The remaining relation has nonzero constant
coefficient $g\in A$, and its evaluation shows $g\in bB$. Consequently
$\pi(Z)\subseteq V(g)$. A nonzero complex polynomial has zero set with
empty Euclidean interior: if it vanishes on a product of open disks,
view it as a polynomial in the last variable, use the finite-root
bound, and induct on the number of variables to show every coefficient
is zero. Since $\pi$ is open, the image of every nonempty Euclidean
open subset of $Y$ meets the complement of $V(g)$. The subset therefore
meets $Y\setminus Z$, proving density. In dimension zero the coordinate
domain is finite over $\mathbb C$, hence equals $\mathbb C$, and the
same conclusion is immediate.

The reverse closure inclusion follows because polynomial zero sets are
Euclidean closed. The statement is
\cite[Aside 9.8]{milne-algebraic-geometry}; the finite-normalization and
root-bound argument here supplies a proof of the density assertion
stated without proof in that aside.
\end{proof}

